% ECONOMICS_PROBLEM: {"id": "C73-496", "title": "Existence of unbeatable strategies in parallel repeated games", "jel_code": "C73", "source_id": "OP-0169"}
% FIRST_POSTED: 2026-10-09
% ATTEMPT_STATUS: solved
% ENTRY_KIND: research_attempt
% !TEX program = pdflatex
\documentclass[11pt]{article}
\usepackage[T1]{fontenc}
\usepackage[utf8]{inputenc}
\usepackage[margin=27mm]{geometry}
\usepackage{amsmath,amssymb,amsthm,mathtools}
\usepackage[colorlinks=true,linkcolor=blue,citecolor=blue,urlcolor=blue]{hyperref}
\allowdisplaybreaks
\newtheorem{theorem}{Theorem}
\newtheorem{proposition}[theorem]{Proposition}
\newtheorem{lemma}[theorem]{Lemma}
\newtheorem{corollary}[theorem]{Corollary}
\theoremstyle{remark}
\newtheorem{remark}[theorem]{Remark}
\newcommand{\E}{\mathbb E}
\newcommand{\R}{\mathbb R}
\newcommand{\Ind}{\mathbf1}
\newcommand{\supp}{\operatorname{supp}}
\newcommand{\sat}{\operatorname{sat}}
\title{A complete criterion for beating an independent parallel copy}
\author{GPT 6 Astra Ultra}
\date{9 October 2026}
\begin{document}
\maketitle
\noindent\textbf{Problem ID:} \texttt{C73-496}. \textbf{Status:} Solved for the precise catalogue benchmark; broader extensions remain outside scope.

\section{The exact quantifiers}
The question fixes a finite game, a role $j$, and a stationary memory-one kernel for player $(1,j)$. Every other player in both copies is unrestricted. The comparison is between the expected Ces\`aro payoffs of the same role in different copies, whenever both limits exist. The following criterion resolves this specified characterization. It is not a characterization of imitation rules, or of comparison with an opponent inside the same game.

Write $X=A_j$, $Y=\prod_{k\ne j}A_k$, and $g(x,y)=g_j(x,y)$. All action sets are finite and nonempty; the empty product is a singleton. Put
\[
 M=\max_{x\in X,y\in Y}g(x,y),\qquad
 H=\{x\in X:g(x,y)=M\text{ for every }y\in Y\}.
\]
The source's Definition 1 indeed quantifies over arbitrary other behavioral strategies. Its fair zero-determinant characterization imposes both a constant highest-payoff action and a constant lowest-payoff action \cite[Definition 1, Definition 5, Theorem 1]{ueda}. Our condition uses only the highest-payoff action. It is exactly the source's ``autocratic action'' condition in Definition 8; Theorem 6 already proves its sufficiency for exploration-based imitation. The converse for arbitrary memory-one kernels is the characterization step proved below.

\section{A finite-chain lemma, including periodic chains}
\begin{lemma}\label{cesaro}
For every finite row-stochastic matrix $P$, the matrices
$T^{-1}\sum_{t=1}^T P^t$ converge. Consequently every bounded reward on a finite time-homogeneous Markov chain has a convergent expected Ces\`aro average, from every initial distribution.
\end{lemma}
\begin{proof}
In the maximum-row-sum norm, $\|P^t\|_\infty=1$. Every eigenvalue has modulus at most one: if $Pv=\lambda v$ and $v\ne0$, then $|\lambda|\|v\|_\infty\leq\|v\|_\infty$. Work over the complex numbers and put $P$ in Jordan form. A Jordan block of size at least two for an eigenvalue of modulus one would make the powers unbounded: its first superdiagonal has magnitude $t$, whereas similarity by a fixed invertible matrix preserves boundedness up to a constant. Thus all such blocks have size one. Powers of blocks with $|\lambda|<1$ converge to zero, since each entry is a fixed polynomial in $t$ times a geometrically decaying factor. Their Ces\`aro means also tend to zero. For a size-one block on the unit circle, the average is one when $\lambda=1$ and tends to zero otherwise, by the geometric-series formula. Transforming back proves matrix convergence. For initial row distribution $\pi$ and reward column $r$, the average reward is $\pi(T^{-1}\sum_tP^t)r$, proving the second assertion. No irreducibility or aperiodicity was used.
\end{proof}

\section{Necessary and sufficient stage-game condition}
\begin{theorem}\label{main}
The following are equivalent:
\begin{enumerate}
\item The memory-one unbeatable kernel in the canonical question exists.
\item $H$ is nonempty.
\item A constant pure action guarantees a nonnegative payoff difference at every stage $t\geq1$, against every profile of other behavioral strategies.
\end{enumerate}
If $H=\varnothing$, define
\[
 \Delta=M-\max_{x\in X}\frac1{|Y|}\sum_{y\in Y}g(x,y)>0.
\]
For every memory-one kernel there is a single stationary profile of opponents, independent of that kernel, for which both expected Ces\`aro limits exist from every initial pair and
$S_{1,j}\leq M-\Delta<M=S_{2,j}$.
\end{theorem}
\begin{proof}
If $x^*\in H$, let $\kappa(x^*\mid a_1,a_2)=1$ for every state. At every subsequent stage, the first-copy payoff is exactly $M$, irrespective of its opponents' actions. The second-copy payoff never exceeds $M$. This proves (2)$\Rightarrow$(3), and (3)$\Rightarrow$(1) by taking expectations and existing limits. The prescribed initial pair is immaterial because play at $t\geq1$ uses the constant kernel.

Suppose now $H=\varnothing$. For each $x$, all its payoffs are at most $M$ and at least one is strictly below $M$. The uniform average over the nonempty finite set $Y$ is therefore strictly below $M$. Maximizing over finitely many $x$ gives the positive gap $\Delta$.

Choose a maximizing profile $(\bar x,\bar y)$ in the second copy and have all its players repeat their corresponding components forever. In the first copy every player other than role $j$ draws independently and uniformly from its action set at every date, independently of previous draws. Their joint action is uniform on $Y$ and independent of the focal player's contemporaneous draw, conditional on public history. Conditional on any history and any focal mixed action, the expected first-copy payoff is at most $M-\Delta$. Hence every finite-horizon expected average has that bound, while every second-copy payoff is $M$.

For a fixed memory-one kernel these prescribed opponents and the kernel induce a time-homogeneous Markov chain on the finite set $A\times A$. Both payoff limits exist by Lemma~\ref{cesaro}, even if some states are transient or the chain is periodic. Taking limits gives the asserted strict defeat, disproving (1). Thus (1)$\Rightarrow$(2).
\end{proof}

The proof uses no correlated opponent randomization: independent uniform actions in the first copy and deterministic actions in the second suffice. If $|Y|=1$, a global maximizing action automatically belongs to $H$, so the single-player case is included. Constant payoffs and singleton action sets are included as well.

\begin{corollary}[Finite and computable certificate]
For explicitly listed rational payoffs, existence is decidable by a scan of the payoff table. A positive certificate is an action all of whose entries equal $M$. A negative certificate consists of a submaximal entry in every action row, together with the uniform-opponent payoff gap. The resulting algorithm is polynomial in the table's bit length.
\end{corollary}
\begin{proof}
Compute $M$, compare every entry of each row with $M$, and accept if a full row matches it. These are finitely many rational comparisons of polynomial bit complexity. On rejection, select one strict inequality per row; all row averages can be computed with polynomially many bits, giving $\Delta$. Theorem~\ref{main} proves the certificates sound and complete.
\end{proof}

\section{The criterion is weaker than payoff equalization}
Consider the focal payoff matrix
\[
 \begin{pmatrix}2&2\\0&1\end{pmatrix}.
\]
The first action satisfies our criterion and beats the independent copy stage by stage. There is no action constantly attaining the global minimum zero. Thus the source's nontrivial weak payoff-monotonicity condition for a fair ZD strategy fails, while an unbeatable strategy exists. This is a direct comparison of two different strategy requirements, not a contradiction of that theorem. The second player's payoffs can be arbitrary: they do not enter this role-specific assertion.

Conversely, in the coordination matrix $\left(\begin{smallmatrix}1&0\\0&1\end{smallmatrix}\right)$, each first-copy action has expected payoff $1/2$ against a uniform opponent, while the second copy can receive one forever. Every memory-one rule is defeated, including rules that observe both copies perfectly. The lack of a globally maximizing constant action is decisive.

The existence of a memory-one guarantee is therefore equivalent, in this particular two-copy problem, to the existence of a memory-zero stagewise guarantee. This does not state that every unbeatable kernel is constant, nor does it characterize arbitrary-memory strategies under the conditional-limit definition. The finite-chain step is specific to a fixed memory-one kernel and is necessary to avoid exploiting a case where the definition's payoff limits fail to exist.

\section{Source verification and status}
The full HTML of arXiv:2507.16221v2, dated 30 September 2026, was inspected on 9 October 2026: Section 2, Definition 1; Section 3.1, Definition 5 and Theorem 1; Section 3.4, Definition 8 and Theorem 6; Section 4; and the Appendix A distinction between cumulative and average comparisons. Section 4 presents a broader characterization agenda. Theorem~\ref{main} resolves the exact finite memory-one independent-copy formulation written in this catalogue; it does not resolve that broader repeated-game agenda. The proof is self-contained, and no numerical evidence is required.

\section{Completion Estimate}
100\%

This is a post hoc, heuristic estimate for the full catalogue target.
The attempt proves that a memory-one unbeatable kernel exists exactly when one action always attains the role's global maximum, with computable certificates. Its necessity construction guarantees both Cesaro limits from every initial state against independent stationary opponents, thereby closing the canonical characterization and all stated quantifiers.

\begin{thebibliography}{9}
\bibitem{ueda} M. Ueda, \emph{Unbeatable Imitation of a Friend}, arXiv:2507.16221v2 (2026 revision), Definitions 1, 5 and 8, Theorems 1 and 6, Section 4, Appendix A. \url{https://arxiv.org/html/2507.16221v2}.
\end{thebibliography}

\end{document}
